% =====================================================================
%  Literature Review  --  standalone working document
%  "Viewpoint Invariance Is a Theorem, Not a Fit:
%   An SE(3)-Equivariant Recurrence for Rehabilitation Exercise Assessment"
%
%  Self-contained. Compiles with:  pdflatex literature_review && pdflatex literature_review
%  Deliberately uses the plain `article` class (NOT IEEEtran/wacv) so it builds
%  fast and independently of the main manuscript's template.
%
%  Every \cite key below is grounded in a real reference in the bibliography at
%  the end of this file; nothing is invented.
% =====================================================================
\documentclass[11pt]{article}

\usepackage[margin=1in]{geometry}
\usepackage{amsmath,amssymb}
\usepackage{booktabs}
\usepackage{array}
\usepackage{enumitem}
\usepackage[table]{xcolor}
\usepackage[colorlinks,allcolors=blue!60!black]{hyperref}

\title{\textbf{Literature Review}\\[4pt]
\large Viewpoint Invariance as a Structural Guarantee for\\
Rehabilitation Exercise Assessment from Skeleton Sequences}
\author{}
\date{}

\newcommand{\gap}[1]{\textbf{\textcolor{red!60!black}{Gap:} } #1}

\begin{document}
\maketitle

\begin{abstract}
\noindent This document surveys the work our paper builds on and departs from. We
organise it along the six axes the paper touches --- (i) rehabilitation exercise
assessment from skeletons, (ii) skeleton-based action recognition, (iii)
group-equivariant and invariant architectures, (iv) learned versus structural
viewpoint invariance, (v) continuous-time models for irregular sampling, and
(vi) edge deployment of structured models --- and after each we state the
specific gap that motivates a contribution. A synthesis table (\S\ref{sec:synth})
maps every gap to the contribution that closes it. The through-line is a single
distinction the existing literature does not make operational: a symmetry that
holds \emph{in expectation over a training distribution} (augmentation) versus one
that holds \emph{identically for every input} (architecture). Our thesis is that
for a clinical screening instrument only the latter is admissible, and that it is
provable, certifiable, and deployable.
\end{abstract}

% =====================================================================
\section{Scope and Problem Statement}
\label{sec:scope}

Automated assessment of physical rehabilitation exercises from skeleton sequences
promises objective, remote, continuous monitoring of therapy~\cite{capecci2019kimore}.
A model receives a $3$D skeleton trajectory --- typically from a consumer depth
sensor --- and regresses a clinical quality score. The dominant evaluation regime
ranks architectures by a single aggregate error against a physician's score on
clean, frontal, single-sensor recordings~\cite{karlov2024,abedi2023,kuang2026,pct2026}.

Our paper argues that this regime, not the architectures, is the limiting factor:
it rewards capacity and penalises nothing a deployed home device actually
encounters --- a camera against a different wall, a tracking node that freezes, a
dropped frame. The review below is therefore structured to expose two recurring
blind spots in the prior art: (a) evaluation is uniformly \emph{in-distribution},
and (b) no prior work on this task reports a \emph{mean-predictor floor}, without
which a robustness curve cannot be interpreted (a model that ignores its input has
zero degradation and would ``win'' every robustness plot).

% =====================================================================
\section{Rehabilitation Exercise Assessment from Skeletons}
\label{sec:rehab}

\paragraph{Benchmarks.} KIMORE~\cite{capecci2019kimore} is the standard corpus:
subjects perform five trunk and hip exercises, each scored $0$--$50$ by a
clinician, with paired healthy and clinical populations. UI-PRMD~\cite{vakanski2018uiprmd}
provides ten movements with marker-based ground truth; IRDS (the
IntelliRehab dataset) contributes gesture-level rehabilitation recordings.
Ismail-Fawaz et al.'s Rehab-Pile~\cite{rehabpile2026} aggregates KIMORE, UI-PRMD and
IRDS under cross-subject splits and reports error rather than rank correlation ---
its protocol discipline is the closest antecedent of our own protocol audit.

\paragraph{Methods.} Reported approaches span a wide architectural range:
ST-GCN variants trained with supervised contrastive objectives using hard and soft
negatives~\cite{karlov2024}; cross-modal video-to-joints augmentation to enlarge
scarce training data~\cite{abedi2023}; dual-stream spatio-temporal graph networks
with motion-aware grouping~\cite{kuang2026}; and, most relevant to us, point-cloud
transformers that treat the skeleton as an unordered set of $3$D points with axial
self-attention~\cite{pct2026}. The last is the baseline we position against
directly.

\paragraph{What the field measures, and what it does not.} Two properties recur.
First, evaluation is uniformly in-distribution: same sensor, same camera placement,
no simulated hardware failure. Second, protocols differ enough that headline
numbers are not comparable --- reported MADs on the most common single-exercise
slice mix optimistic epoch-selection rules with subject-disjointness assumptions
that do not hold uniformly.

\gap{No prior rehab-assessment work (a)~stress-tests off the frontal, clean,
single-sensor manifold, or (b)~reports a mean-predictor floor against which a
degradation curve is interpretable. Protocol incomparability is acknowledged
anecdotally but never isolated quantitatively.}

% =====================================================================
\section{Skeleton-Based Action Recognition}
\label{sec:action}

The methodological backbone of skeleton scoring is the action-recognition
literature. Spatial-temporal graph convolutional networks (ST-GCN)~\cite{yan2018stgcn}
established the graph-over-joints, convolution-over-time paradigm that most
rehabilitation models inherit, and remain the standard baseline for cross-view and
cross-subject generalisation on NTU~RGB+D~\cite{shahroudy2016ntu}. NTU~RGB+D's
\emph{Cross-View} protocol --- train on two physical camera setups, test on a third
displaced by $45^\circ$ --- is the field's canonical \emph{real} (non-synthetic)
viewpoint-generalisation benchmark, and we use it precisely to show our viewpoint
guarantee is not an artifact of synthetic test-time rotation.

\gap{Cross-view robustness in action recognition is pursued almost exclusively by
augmentation or view canonicalisation, i.e.\ as a statistical property of a trained
model over a view distribution --- never as an exact, per-prediction architectural
identity. Occlusion/tracking failure is treated as a nuisance to be masked away,
not as a discriminating experiment between otherwise-equivalent architectures.}

% =====================================================================
\section{Equivariant and Invariant Neural Architectures}
\label{sec:equiv}

Group-equivariant networks impose symmetry as a \emph{constraint on the hypothesis
class} rather than a target of fitting~\cite{cohen2016gconv}. For $3$D geometry the
steerable family --- tensor-field networks~\cite{thomas2018tfn},
$SE(3)$-Transformers~\cite{fuchs2020se3transformer}, and the \texttt{e3nn}
framework~\cite{geiger2022e3nn} --- carries features in irreducible representations
of $O(3)$ and constrains every layer to intertwine with the Wigner-$D$ action,
guaranteeing equivariance \emph{for arbitrary weights} (trained, untrained, or
adversarial). Lighter constructions obtain $E(n)$-equivariance without spherical
harmonics~\cite{satorras2021egnn}, or lift a network to $SO(3)$-equivariance by
replacing scalar neurons with vector-valued ones~\cite{deng2021vectorneurons}. An
orthogonal line makes a non-equivariant backbone invariant by canonicalisation or
frame averaging~\cite{puny2022frameaveraging} --- exact, but at a forward-pass or
learned-canonicaliser cost.

These architectures are mature in molecular modelling and point-cloud vision but
have not been brought to bear on rehabilitation scoring as an \emph{exact viewpoint
guarantee}. Two further points are, to our knowledge, unstated in this literature:
(a) the natural invariant read-out one forms from such an encoder --- built from
norms, inner products, and bone lengths --- is composed entirely of parity-\emph{even}
features and is therefore silently invariant to a \emph{mirror} of the whole body,
i.e.\ complete for $O(3)$ rather than $SO(3)$, and thus blind to trajectory
handedness until pseudo-scalars are admitted; and (b) whether the steerable
machinery is \emph{necessary} for accuracy is rarely tested against a hand-crafted
invariant baseline of equal guarantee.

\gap{Equivariant architectures give an exact viewpoint guarantee ``for free'' but
(i)~this has not been certified numerically for a rehabilitation read-out, (ii)~its
parity/chirality completeness is unexamined, and (iii)~its necessity over a simple
hand-crafted invariant model has not been isolated by a discriminating experiment.}

% =====================================================================
\section{Learned versus Structural Viewpoint Invariance}
\label{sec:learned}

The default remedy for viewpoint variation is training-time augmentation, and it
``works'' in the sense the field measures: an augmented model's \emph{aggregate}
error becomes approximately flat in camera angle. But augmentation makes prediction
errors \emph{cancel} in the population mean rather than \emph{vanish} per input.
The distinction between a symmetry that holds in expectation over a training
distribution and one that holds identically for every input is, for a clinical
instrument, the whole point --- a screening tool must return the \emph{same patient
the same score from a different room} --- yet it is expressible only if degradation
is reported \emph{per sequence}, not as an aggregate.

\gap{The field reports aggregate accuracy versus camera angle, which a
rotation-augmented model can flatten. It does not report per-sequence score shift,
so a model whose errors merely cancel is indistinguishable from one that is
genuinely invariant.}

% =====================================================================
\section{Continuous-Time Models and Irregular Sampling}
\label{sec:cde}

Depth-sensor streams are not uniformly sampled --- genuine multi-frame dropouts and
jittered inter-arrival times are the norm. The standard treatment resamples onto a
uniform grid; the standard alternative is a continuous-time model consuming arrival
times natively: Neural ODEs~\cite{chen2018neuralode}, latent ODEs and ODE-RNNs for
irregular series~\cite{rubanova2019odernn}, Neural CDEs driven by an interpolant of
the observations~\cite{kidger2020neuralcde}, and recurrences with explicit decay
between observations (GRU-D)~\cite{che2018grud}, all built on gated recurrent
units~\cite{cho2014gru}. The widely assumed benefit is that native consumption of
irregular timestamps must beat resampling.

Two subtleties in this literature are load-bearing for us. First, adaptive
solvers such as Dormand--Prince~\cite{dormand1980prince} accept or reject steps on a
\emph{scalar} reduction of the embedded local error; if that scalar is not
group-invariant, a rotated and an unrotated integration receive \emph{different
step grids} and silently de-align --- even when the vector field is exactly
equivariant and the representation orthogonal. Second, irregularly sampled data have
no single Nyquist frequency and do not alias, which is the basis of the
Lomb--Scargle construction~\cite{vanderplas2018lombscargle}: a grid-free estimator
can recover a high-frequency tone that a uniform grid at the same \emph{average}
rate folds onto itself irrecoverably.

\gap{The advantage of native irregular-time consumption is assumed, not conditioned
on the signal's spectral content; and the interaction between adaptive-solver error
control and group equivariance is undocumented, so a ``certified equivariant''
continuous-time model can still silently break invariance through step-grid
divergence.}

% =====================================================================
\section{Edge Deployment of Structured Models}
\label{sec:edge}

A guarantee proved in exact ($64$-bit) arithmetic on a full recording is not the
same as a guarantee that survives on the hardware a home device runs: reduced
precision (fp16/bf16/int8) and a real-time latency budget. Quantisation is not an
orthogonal transformation, so it does not commute with the group action, and a
high-accuracy skeleton classifier that is bidirectional and pools over the whole
clip cannot emit a score until the last frame --- it has no honest latency.

\gap{The precision budget of an equivariance guarantee (where, exactly, the theorem
degrades from fp64 to int8, and whether it is weight- or activation-quantisation
that breaks it) is unmeasured; and viewpoint-robust skeleton scorers are not
co-designed for strictly causal, frame-by-frame real-time operation.}

% =====================================================================
\section{Synthesis: Gaps and How This Work Closes Them}
\label{sec:synth}

Table~\ref{tab:synth} collects the gaps above and maps each to the contribution
that addresses it. The unifying claim is that clean-accuracy leaderboards cannot
separate architectures on this task, and that the decisive variable is
\emph{structural behaviour off the clean, frontal manifold}.

\begin{table}[h]
\centering
\small
\renewcommand{\arraystretch}{1.35}
\begin{tabular}{p{0.30\linewidth} p{0.62\linewidth}}
\toprule
\textbf{Gap in prior work} & \textbf{How this work closes it} \\
\midrule
No mean-predictor floor; protocol incomparability unquantified
  & Protocol audit isolates optimistic epoch selection as a $16.8\%$ effect on the
    standard KIMORE slice, shows that slice carries no subject-level signal, fixes a
    pooled subject-disjoint protocol, and gates every claim on clearing the floor. \\
Exact viewpoint invariance never certified for a rehab read-out
  & An $SE(3)$-invariant read-out over a steerable encoder~\cite{geiger2022e3nn},
    audited by eight numerical gates (E1--E8): orthogonality $\le10^{-12}$, violation
    at machine roundoff with a \emph{flat} slope in rotation angle, an $8\times10^{8}$
    fp32/fp64 ratio; cross-viewpoint degradation $9\times10^{-6}$~MAD vs.\ the
    baseline's $9.42$. \\
Viewpoint robustness only shown on synthetic rotation
  & Confirmed on NTU~RGB+D~$60$ Cross-View~\cite{shahroudy2016ntu,yan2018stgcn}, a
    \emph{real} $45^\circ$ cross-camera benchmark, at exact ($5.8\times10^{-9}$, fp64)
    cross-view logit invariance without rotation augmentation. \\
Augmentation's cost reported only in aggregate
  & The rotation-augmented baseline is granted its remedy and its price measured: a
    flat $+1.0$~MAD penalty \emph{and} a $3.03$~MAD per-sequence score shift that the
    aggregate hides --- a flat accuracy curve is shown to \emph{not} be invariance. \\
Parity/chirality completeness of invariant cuts unexamined
  & The natural invariant cut is exhibited (gate E6) to be silently $O(3)$-invariant
    and blind to trajectory handedness; pseudo-scalars restore $SO(3)$ at no cost to
    the viewpoint theorem (E8). \\
Necessity of steerable machinery not isolated
  & A GRU on hand-crafted $SO(3)$ invariants matches accuracy \emph{and} the exact
    guarantee at a third of the parameters; the steerable encoder is then justified by
    a \emph{measured} sensor-node-failure experiment, not rhetoric. \\
Irregular-sampling advantage assumed, not conditioned
  & Stated as a bandwidth condition and verified in both directions: a spectral census
    shows resampling destroys only $2.2\%$ of KIMORE's positional energy (a denoiser),
    predicting the null; a $5$~Hz tremor placed in the destroyed band is recovered at
    $r=1.000$ under $70\%$ drops via a Lomb--Scargle-style estimator~\cite{vanderplas2018lombscargle}. \\
Solver error control vs.\ equivariance undocumented
  & The per-component RMS norm of a Dormand--Prince controller~\cite{dormand1980prince}
    is proved non-invariant; a per-irrep-normalised norm drives step-grid divergence to
    identically zero. The Neural CDE~\cite{kidger2020neuralcde} is certified to the same
    standard yet \emph{fails} the floor --- a clean demonstration that a structural
    guarantee is worthless without measured signal. \\
Precision budget and real-time latency of the guarantee unmeasured
  & The equivariance floor is measured fp64$\to$int8; weight-quantisation preserves the
    theorem exactly while only activation-rounding breaks it; the read-out streams
    causally at $8.9$~ms/frame ($112$~fps) with an honest time-to-first-score. \\
\bottomrule
\end{tabular}
\caption{Gap-to-contribution map. Each row is a shortfall identified in
\S\ref{sec:rehab}--\S\ref{sec:edge} and the measured result that addresses it.}
\label{tab:synth}
\end{table}

% =====================================================================
\section{Positioning of Our Novelty}
\label{sec:novelty}

Relative to the literature above, the paper's novelty is not a new equivariant
layer or a new benchmark score. It is a \emph{reframing of what should be measured}
on this task, plus the structural machinery and certification to act on it:

\begin{enumerate}[leftmargin=1.4em]
\item \textbf{Assessment as an invariant read-out of an equivariant representation.}
We formulate rehabilitation scoring so that invariance of the predicted score under
any rigid motion of the camera is a \emph{theorem about the architecture}, true for
arbitrary weights, rather than a behaviour induced by augmentation --- and we
\emph{certify} it numerically (E1--E8) rather than assert it.

\item \textbf{The measurement discipline the task was missing.} A mean-predictor
floor, a quantified nondeterminism floor ($\approx0.33$~MAD) below which no accuracy
gap is real, and per-sequence degradation reported alongside aggregate accuracy
everywhere. This turns ``robustness'' from an uninterpretable plot into a gated
claim.

\item \textbf{Two negative results, one of them derived.} We report --- rather than
hide --- that native irregular-time consumption buys nothing on KIMORE, and we
\emph{predict} the null from a spectral census before running the experiment; and
that restoring chirality is principled but unrewarded on a corpus of healthy-form
exercises. Honest scope boundaries, with the corpus properties that would reward
each capability named explicitly.

\item \textbf{A discriminating experiment, not an appeal to elegance.} Because a
hand-crafted invariant GRU matches the steerable encoder on both accuracy and the
viewpoint guarantee, we justify the steerable machinery by the one axis on which
they measurably differ --- graceful degradation under a frozen tracking node ---
and we record, against ourselves, that the point-cloud transformer is in fact the
most node-robust arm, so the defensible claim is a Pareto one.

\item \textbf{The guarantee as a deployable systems property.} A structural
quantisation audit (weight- vs.\ activation-rounding) and a strictly causal
streaming read-out show the theorem is co-designed for the edge, not merely an
fp64 identity.
\end{enumerate}

\noindent In one sentence: prior work optimises clean, frontal, aggregate accuracy;
we show that on this task that quantity cannot separate architectures, and that the
decisive, clinically relevant variable --- a per-prediction viewpoint guarantee that
survives real cameras, a broken sensor, reduced precision, and real-time latency ---
is \emph{structural}, provable, and measurable.

% =====================================================================
\begin{thebibliography}{00}

\bibitem{capecci2019kimore}
M.~Capecci, M.~G.~Ceravolo, F.~Ferracuti, S.~Iarlori, A.~Monteriu, L.~Romeo, and F.~Verdini,
``The KIMORE dataset: Kinematic assessment of movement and clinical scores for remote monitoring of physical rehabilitation,''
\textit{IEEE Trans. Neural Syst. Rehabil. Eng.}, vol.~27, no.~7, pp.~1436--1448, 2019.

\bibitem{pct2026}
``A point cloud transformer for remote monitoring and automated assessment of physical rehabilitation exercises,''
arXiv:2606.30309, 2026.

\bibitem{karlov2024}
M.~Karlov, A.~Abedi, and S.~S.~Khan,
``Rehabilitation exercise quality assessment through supervised contrastive learning with hard and soft negatives,''
\textit{Med. Biol. Eng. Comput.}, 2024.

\bibitem{abedi2023}
A.~Abedi, M.~Malmirian, and S.~S.~Khan,
``Cross-modal video to body-joints augmentation for rehabilitation exercise quality assessment,''
arXiv:2306.09546, 2023.

\bibitem{kuang2026}
Z.~Kuang, Y.~Yin, Y.~Yang, J.~Zhao, and L.~Sun,
``Dual-stream STGCN with motion-aware grouping for rehabilitation action quality assessment,''
\textit{Sensors}, vol.~26, no.~1, p.~287, 2026.

\bibitem{rehabpile2026}
A.~Ismail-Fawaz et al.,
``Rehab-Pile: A cross-subject rehabilitation benchmark aggregating KIMORE, UI-PRMD and IRDS,''
arXiv:2507.21018, IEEE FG 2026.

\bibitem{vakanski2018uiprmd}
A.~Vakanski, H.-P.~Jun, D.~Paul, and R.~Baker,
``A data set of human body movements for physical rehabilitation exercises (UI-PRMD),''
\textit{Data}, vol.~3, no.~1, p.~2, 2018.

\bibitem{yan2018stgcn}
S.~Yan, Y.~Xiong, and D.~Lin,
``Spatial temporal graph convolutional networks for skeleton-based action recognition,''
in \textit{Proc. AAAI}, 2018.

\bibitem{shahroudy2016ntu}
A.~Shahroudy, J.~Liu, T.-T.~Ng, and G.~Wang,
``NTU RGB+D: A large scale dataset for 3D human activity analysis,''
in \textit{Proc. IEEE CVPR}, 2016.

\bibitem{cohen2016gconv}
T.~S.~Cohen and M.~Welling,
``Group equivariant convolutional networks,''
in \textit{Proc. ICML}, 2016.

\bibitem{thomas2018tfn}
N.~Thomas, T.~Smidt, S.~Kearnes, L.~Yang, L.~Li, K.~Kohlhoff, and P.~Riley,
``Tensor field networks: Rotation- and translation-equivariant neural networks for 3D point clouds,''
arXiv:1802.08219, 2018.

\bibitem{fuchs2020se3transformer}
F.~B.~Fuchs, D.~E.~Worrall, V.~Fischer, and M.~Welling,
``SE(3)-Transformers: 3D roto-translation equivariant attention networks,''
in \textit{Proc. NeurIPS}, 2020.

\bibitem{geiger2022e3nn}
M.~Geiger and T.~Smidt,
``e3nn: Euclidean neural networks,''
arXiv:2207.09453, 2022.

\bibitem{satorras2021egnn}
V.~G.~Satorras, E.~Hoogeboom, and M.~Welling,
``E(n) equivariant graph neural networks,''
in \textit{Proc. ICML}, 2021.

\bibitem{deng2021vectorneurons}
C.~Deng, O.~Litany, Y.~Duan, A.~Poulenard, A.~Tagliasacchi, and L.~Guibas,
``Vector neurons: A general framework for SO(3)-equivariant networks,''
in \textit{Proc. ICCV}, 2021.

\bibitem{puny2022frameaveraging}
O.~Puny, M.~Atzmon, H.~Ben-Hamu, I.~Misra, A.~Grover, E.~J.~Smith, and Y.~Lipman,
``Frame averaging for invariant and equivariant network design,''
in \textit{Proc. ICLR}, 2022.

\bibitem{chen2018neuralode}
R.~T.~Q.~Chen, Y.~Rubanova, J.~Bettencourt, and D.~Duvenaud,
``Neural ordinary differential equations,''
in \textit{Proc. NeurIPS}, 2018.

\bibitem{rubanova2019odernn}
Y.~Rubanova, R.~T.~Q.~Chen, and D.~Duvenaud,
``Latent ODEs for irregularly-sampled time series,''
in \textit{Proc. NeurIPS}, 2019.

\bibitem{kidger2020neuralcde}
P.~Kidger, J.~Morrill, J.~Foster, and T.~Lyons,
``Neural controlled differential equations for irregular time series,''
in \textit{Proc. NeurIPS}, 2020.

\bibitem{che2018grud}
Z.~Che, S.~Purushotham, K.~Cho, D.~Sontag, and Y.~Liu,
``Recurrent neural networks for multivariate time series with missing values,''
\textit{Sci. Rep.}, vol.~8, no.~1, p.~6085, 2018.

\bibitem{cho2014gru}
K.~Cho, B.~van Merri\"{e}nboer, C.~Gulcehre, D.~Bahdanau, F.~Bougares, H.~Schwenk, and Y.~Bengio,
``Learning phrase representations using RNN encoder--decoder for statistical machine translation,''
in \textit{Proc. EMNLP}, 2014.

\bibitem{dormand1980prince}
J.~R.~Dormand and P.~J.~Prince,
``A family of embedded Runge--Kutta formulae,''
\textit{J. Comput. Appl. Math.}, vol.~6, no.~1, pp.~19--26, 1980.

\bibitem{vanderplas2018lombscargle}
J.~T.~VanderPlas,
``Understanding the Lomb--Scargle periodogram,''
\textit{Astrophys. J. Suppl. Ser.}, vol.~236, no.~1, p.~16, 2018.

\end{thebibliography}

\end{document}
